\documentclass[pract]{SCWorks}
% Тип обучения (одно из значений):
%    bachelor   - бакалавриат (по умолчанию)
%    spec       - специальность
%    master     - магистратура
% Форма обучения (одно из значений):
%    och        - очное (по умолчанию)
%    zaoch      - заочное
% Тип работы (одно из значений):
%    coursework - курсовая работа (по умолчанию)
%    referat    - реферат
%  * otchet     - универсальный отчет
%  * nirjournal - журнал НИР
%  * digital    - итоговая работа для цифровой кафедры
%    diploma    - дипломная работа
%    pract      - отчет о научно-исследовательской работе
%    autoref    - автореферат выпускной работы
%    assignment - задание на выпускную квалификационную работу
%    review     - отзыв руководителя
%    critique   - рецензия на выпускную работу

% * Добавлены вручную. За вопросами к @mchernigin
\usepackage{preamble}

\begin{document}

% Кафедра (в родительном падеже)
\chair{математической кибернетики и компьютерных наук}

% Тема работы
\title{Модули в языке Python}

% Курс
\course{3}

% Группа
\group{351}

% Факультет (в родительном падеже) (по умолчанию "факультета КНиИТ")
% \department{факультета КНиИТ}

% Специальность/направление код - наименование
% \napravlenie{02.03.02 "--- Фундаментальная информатика и информационные технологии}
% \napravlenie{02.03.01 "--- Математическое обеспечение и администрирование информационных систем}
% \napravlenie{09.03.01 "--- Информатика и вычислительная техника}
\napravlenie{09.03.04 "--- Программная инженерия}
% \napravlenie{10.05.01 "--- Компьютерная безопасность}

% Для студентки. Для работы студента следующая команда не нужна.
\studenttitle{студентки}

% Фамилия, имя, отчество в родительном падеже
\author{Винник Маргариты Дмитриевны}

% Заведующий кафедрой 
\chtitle{доцент, к.\,ф.-м.\,н.}
\chname{С.\,В.\,Миронов}

% Руководитель ДПП ПП для цифровой кафедры (перекрывает заведующего кафедры)
% \chpretitle{
%     заведующий кафедрой математических основ информатики и олимпиадного\\
%     программирования на базе МАОУ <<Ф"=Т лицей №1>>
% }
% \chtitle{г. Саратов, к.\,ф.-м.\,н., доцент}
% \chname{Кондратова\, Ю.\,Н.}

% Научный руководитель (для реферата преподаватель проверяющий работу)
\satitle{доцент, к.\,ф.-м.\,н.} %должность, степень, звание
\saname{М.\,И.\,Сафрончик}

% Руководитель практики от организации (руководитель для цифровой кафедры)
\patitle{доцент, к.\,ф.-м.\,н.}
\paname{М.\,И.\,Сафрончик}

% Руководитель НИР
\nirtitle{доцент, к.\,п.\,н.} % степень, звание
\nirname{В.\,А.\,Векслер}

% Семестр (только для практики, для остальных типов работ не используется)
\term{6}

% Наименование практики (только для практики, для остальных типов работ не
% используется)
\practtype{учебная}

% Продолжительность практики (количество недель) (только для практики, для
% остальных типов работ не используется)
\duration{2}

% Даты начала и окончания практики (только для практики, для остальных типов
% работ не используется)
\practStart{08.05.2026}
\practFinish{26.05.2026}

% Год выполнения отчета
\date{2026}

\maketitle

% Включение нумерации рисунков, формул и таблиц по разделам (по умолчанию -
% нумерация сквозная) (допускается оба вида нумерации)
\secNumbering

\tableofcontents

% Раздел "Обозначения и сокращения". Может отсутствовать в работе
% \abbreviations
% \begin{description}
%     \item ... "--- ...
%     \item ... "--- ...
% \end{description}

% Раздел "Определения". Может отсутствовать в работе
% \definitions

% Раздел "Определения, обозначения и сокращения". Может отсутствовать в работе.
% Если присутствует, то заменяет собой разделы "Обозначения и сокращения" и
% "Определения"
% \defabbr

\intro

% После введения — серии \section, \subsection и т.д.

Сам по себе «чистый» Python, безусловно, удобен, но для сложных численных расчётов, анализа данных (Data Science) и 
машинного обучения (Machine learning) его встроенных возможностей недостаточно. Ключевая причина, почему Python стал 
лидером в области анализа данных и искусственного интеллекта — это его способность легко интегрироваться 
с кодом на других языках.

Ядро интерпретатора CPython (самая распространённая реализация Python) написано на C и C++. Этот фундамент обеспечивает 
эффективное управление памятью, работу с объектами и выполнение базовых операций. 
Многие встроенные функции (print, len) и модули стандартной библиотеки (math, sys) также реализованы 
на C для достижения максимальной производительности.

Однако настоящая мощь раскрывается в сторонних библиотеках. Такие «тяжеловесы», как NumPy для векторных вычислений, 
Pandas для анализа табличных данных, а также фреймворки глубокого обучения PyTorch и TensorFlow, являются модулями 
расширения на C/C++.

Эта же философия расширения лежит в основе библиотек для классического машинного обучения и визуализации. 
Библиотека scikit-learn, построенная поверх NumPy и SciPy, содержит оптимизированные реализации алгоритмов 
на C/C++ (например, для вычисления расстояний, построения деревьев решений), предоставляя при этом единый 
Python-интерфейс. 

Для задач визуализации данных фундаментальная библиотека matplotlib наследует возможности векторной графики 
из C++-библиотек (в частности, Anti-Grain Geometry), а высокоуровневая надстройка seaborn позволяет создавать 
сложные статистические графики минимальным количеством кода, полагаясь на ту же C/C++-основу. 

Таким образом, каждая из ключевых библиотек экосистемы Data Science следует принципу "Python-обёртка над 
высокопроизводительным ядром".

Таким образом, каждая из ключевых библиотек экосистемы Data Science
следует принципу "Python-обёртка над высокопроизводительным ядром".

В этом контексте Python выступает в роли высокоуровневого «клея»: он
предоставляет удобный и гибкий интерфейс для сборки и управления высо-
копроизводительными вычислительными ядрами, написанными на компилиру-
емых языках. Такой подход сочетает простоту разработки со скоростью вы-
полнения кода, что критически важно для задач анализа данных и машинного
обучения.

Важно отметить, что все эти библиотеки спроектированы для бесшов-
ной интеграции друг с другом. NumPy обеспечивает единый формат данных
(ndarray), который понимают Pandas, SciPy, Scikit-learn и Matplotlib. Pandas, в
свою очередь, построен поверх NumPy и добавляет удобную работу с табличны-
ми данными и временными рядами. Scikit-learn принимает на вход как NumPy-
массивы, так и Pandas-DataFrame. Seaborn изначально задуман для работы с
DataFrame и интегрируется с Matplotlib для тонкой настройки графиков. Эта
взаимная совместимость создаёт цельную экосистему, где каждая библиотека
усиливает возможности другой, оставаясь при этом частью единого "языкового
клея"— Python.

\section{Теоритическая часть}

О чём собственно я буду рассказывать:

В этой работе я рассмотрю несколько основных библиотек-инструментов (модулями на самом деле их уже 
давно не называет из-за их объёмности и сложности) для специалиста анализа данных и инженера машинного обучения.
Вот он - так называемый джентельметский набор, который я постепенно рассмотрю подробнее, объясню 
что для чего нужно и как использовать:

\begin{enumerate}
    \item NumPy (Numerical Python);
    \item Pandas;
    \item Scikit-learn;
    \item Matplotlib / Seaborn;
\end{enumerate}

\subsection{Библиотека NumPy}

По классике, начнём с NumPy, но немного издалека, а именно с типа данных, который имеют объекты, 
созданные при помощи этой библиотекой. Тип данных \textbf{ndarray} библиотеки NumPy представляет собой
многомерный массив элементов одного типа. Элементы массива NumPy индексируются кортежем целых 
неотрицательных чисел. 

Массивы NumPy похожи на списки (list) Python или на массивы из встроенного модуля \textbf{array}, которые 
используются для создания плотных массивов данных одного типа. Однако массивы NumPy требуют меньше 
памяти и обычно работают быстрее, поскольку используют оптимизированный, прекомпилированный код на 
языке C.

В то время как объект array в языке Python обеспечивает только эффективное хранение данных в формате 
массива, библиотека NumPy добавляет ещё и возможность выполнения эффективных операций с этими данными 
(будут рассмотрены дальше). 

Ключевая особенность — поддержка поэлементных операций. Это позволяет выполнять основные арифметические 
вычисления сразу над всем массивом, делая код компактным и удобочитаемым.
Также доступна поэлементная операция над двумя массивами одинаковой размерности, в результате 
которой получается другой массив той же размерности, где каждый элемент i, j является результатом вычисления, 
выполненного над элементами i, j двух исходных массивов.

\begin{table}[ht]
\centering
\caption{Сравнение типов данных для хранения массивов в Python}
\begin{tabular}{|p{4.7cm}|p{3.3cm}|p{3.3cm}|p{3.5cm}|}
\hline
\textbf{Характеристика} & \textbf{Встроенный list} & \textbf{Встроенный array из модуля array} & \textbf{ndarray из NumPy} \\ \hline
\textbf{Типы данных} & Любые, разнородные & Однородный (один тип) & Однородный (один тип) \\ \hline
\textbf{Размер} & Динамический & Фиксированный & Фиксированный (при создании) \\ \hline
\textbf{Производительность} & Низкая & Средняя & Высокая (благодаря C/C++ ядру) \\ \hline
\textbf{Функциональность} & Базовая (добавление, удаление) & Очень базовая & Огромная (векторизация, линейная алгебра, статистика) \\ \hline
\textbf{Потребление памяти} & Большое & Маленькое & Маленькое (плотное хранение) \\ \hline
\end{tabular}
\label{tab:comparison}
\end{table}

Прежде чем начать работать с этой библиотекой, её нужно установить, т.к. она не является встроенным 
модулем, как, например, модули math, array и другие стандартные.  
Установка NumPy происходит по команде в терминале (VS Code, PyCharm): \textbf{pip install numpy}. Импортируем пакет 
NumPy под псевдонимом \textbf{np} (это уже в рабочем файле .py): 

\begin{minted}{python}
import numpy as np  
\end{minted}

Дальше в примерах кода буду опускать строку с импортом.

Итак, способы создания массивов NumPy (буду сразу приводить примеры кода и давать пояснения):

\begin{itemize}
    \item Массив NumPy можно создать из данных одного или нескольких списков. Для создания массивов из 
    списков можно воспользоваться функцией \textbf{np.array}.
    \begin{minted}{python}
    my_arr = np.array([1, 4, 2, 5, 3])    
    \end{minted}

    Если мы захотим вывести такой массив, то увидим в консоле следующее: 
    
    \textbf{[1 4 2 5 3]}.
    
    С помощью функции \textbf{type()} можем убедиться, что тип созданного объекта \textbf{ndarray}.
    В консоли увидим: \textbf{<class 'numpy.ndarray'>}.
    
    Сразу рассмотрим пример созддания многомерного массива из нескольких "листов": 
    
    Для каждого сотрудника компании существует список выплат базового оклада за последние 
    три месяца. Чтобы собрать всю информацию о зарплате в одну структуру данных, можно использовать 
    следующий код:

    \begin{minted}{python}
    jeff_salary = [2700,3000,3000]
    nick_salary = [2600,2800,2800]
    tom_salary = [2300,2500,2500]
    base_salary = np.array([jeff_salary, nick_salary, tom_salary])
    print(base_salary)
    \end{minted}

    В консоли это будет выглядеть так:

    [[2700 3000 3000]

    [2600 2800 2800]

    [2300 2500 2500]]

    \item Явное описание многомерного массива (матрицы). Вот один из способов инициализации многомерного
    массива с помощью списка списков:

    \begin{minted}{python}
    my_arr = np.array([range(i, i + 3) for i in [2, 4, 6]])
    print(my_arr)
    \end{minted}

    В консоли это будет выглядеть так:

    [[2 3 4]
    
    [4 5 6]
    
    [6 7 8]]

    \item Создание массивов из нулей или единиц заданной размерности:
    \begin{minted}{python}
    my_arr1 = np.zeros((3, 5), dtype=int)
    my_arr2 = np.ones((7, 3), dtype=float)

    print(my_arr1)
    print(my_arr2)
    \end{minted}

    Выведет двумерный массив целых нулей:

    [[0 0 0 0 0]
    
    [0 0 0 0 0]
    
    [0 0 0 0 0]]

    И массив единиц с плавающей точкой:

    [[1. 1. 1.]
    
    [1. 1. 1.]
    
    [1. 1. 1.]
    
    [1. 1. 1.]
    
    [1. 1. 1.]
    
    [1. 1. 1.]
    
    [1. 1. 1.]]

    Сразу возникает вопрос: "Зачем это нужно и кто это использует?".
    
    Напомню, что привычный список в Python - это, по сути, массив указателей. Каждый элемент списка - 
    это отдельный полноценный объект в памяти (со своей метаинформацией, типом и счетчиком ссылок), и эти 
    объекты могут быть разбросаны по оперативной памяти как угодно.
    Чтобы пройтись по списку, процессору нужно постоянно "прыгать" по памяти, собирая данные. 
    
    Массив в NumPy (ndarray) устроен иначе. Он хранит данные единым, непрерывным блоком в памяти (как массивы в языке C). 
    Кроме того, было уже отмечено, что все элементы в массиве NumPy имеют один и тот же тип данных. 
    Процессор быстро обрабатывает подобные структуры: 
    он может загрузить огромный кусок массива в свой быстрый кэш и обработать его за доли секунды. 
    А еще почти вся математика внутри NumPy написана на высокооптимизированном C, что позволяет обходить 
    медленный интерпретатор Python.

    Маленький вывод: В NumPy выделение памяти под массивы работает как в C, т.е. выделяется непрерывный 
    блок в памяти нужного размера. В сочетании с единым типом данных для всех элементов массива 
    и реализаций математических вычислений на C "под капотом" получается скомпенсировать (обойти) 
    медленный интерпретатор Python и получить быстрые вычисления матриц.
    \item Есть ещё много разнобразных способов создать массив NumPy. Сейчас я их кратко продемонстрирую:

    \begin{minted}{python}
    # Cоздаем массив заданной размерности, заполненный 
    # заданным числом
    arr_this_full = np.full((3, 5), 98.2)

    # Cоздаем массив, заполненный линейной последовательностью,
    # начинающейся с 0 и заканчивающейся 20, с шагом 2
    # (действует подобно встроенной функции range())
    arr_this_arange = np.arange(0, 20,  2)

    # Создаем массив с пятью значениями, располагающимися
    # через равные интервалы в диапазоне между 0 и 1
    arr_this_linspace = np.linspace(0, 1, 5)

    # Создаем массив заданной размерности равномерно
    # распределенных случайных значений от 0 до 1
    arr_this_random = np.random.random((3, 3))

    # Создаем массив заданной размерности
    # нормально распределенных случайных значений
    # со средним 0 и стандартным отклонением 1
    arr_this_random_normal = np.random.normal(0, 1, (3, 3))

    # Создаем массив заданной размерности случайных целых чисел
    # в интервале [0, 10)
    arr_this_randint = np.random.randint(0, 10, (3, 3))

    # Создаем единичную матрицу заданной размерности
    arr_this_ones_matrix =  np.eye(5)

    # Создаем неинициализированный массив из заданного числа
    # целочисленных значений. Значениями будут произвольные
    # данные, случайно оказавшиеся в соответствующих 
    # ячейках памяти
    arr_this_empty = np.empty((4, 3))
    \end{minted}

    Вот как это будет выглядеть в консоли:

    \begin{minted}{python}
    print(arr_this_full)
    \end{minted}

    [[98.2 98.2 98.2 98.2 98.2]
    
    [98.2 98.2 98.2 98.2 98.2]
    
    [98.2 98.2 98.2 98.2 98.2]]

    \begin{minted}{python}
    print(arr_this_arange)
    \end{minted}

    [ 0  2  4  6  8 10 12 14 16 18]

    \begin{minted}{python}
    print(arr_this_linspace)
    \end{minted}

    [0.   0.25 0.5  0.75 1.  ]

    \begin{minted}{python}
    print(arr_this_random)
    \end{minted}

    [[0.85842966 0.19472009 0.53980876]
    
    [0.80552811 0.14153011 0.70063978]
    
    [0.54145708 0.56197047 0.79364311]]

    \begin{minted}{python}
    print(arr_this_random_normal)
    \end{minted}

    [[ 1.3439026   0.97320026 -0.39531868]
    
    [ 0.92342143  0.7090699  -0.42076108]
    
    [-0.21460648 -0.23570425 -0.87608365]]

    \begin{minted}{python}
    print(arr_this_randint)
    \end{minted}
    
    [[0 2 1]
    
    [4 4 9]
    
    [7 8 7]]

    \begin{minted}{python}
    print(arr_this_ones_matrix)
    \end{minted}

    [[1. 0. 0. 0. 0.]
    
    [0. 1. 0. 0. 0.]
    
    [0. 0. 1. 0. 0.]
    
    [0. 0. 0. 1. 0.]
    
    [0. 0. 0. 0. 1.]]

    \begin{minted}{python}
    print(arr_this_empty)
    \end{minted}

    [[6.95198423e-310 6.95198425e-310 6.95198424e-310]
    
    [6.95198423e-310 6.95198423e-310 6.95198426e-310]
    
    [6.95198426e-310 6.95198425e-310 6.95198425e-310]
    
    [6.95198425e-310 6.95198426e-310 6.95198424e-310]]

\end{itemize}

Массивы NumPy содержат значения простых типов, поэтому важно знать и понимать возможности и 
ограничения этих типов, хотя вообще-то их немало и даже предоставлена возможность работать с
комплексными числами.

\begin{longtable}{|p{2.4cm}|p{14cm}|}
\caption{Стандартные типы данных библиотеки NumPy} \label{tab:numpy-types} \\
\hline
\textbf{Тип данных} & \textbf{Описание} \\ \hline
\endfirsthead
\caption*{Продолжение таблицы \ref{tab:numpy-types}} \\
\hline
\textbf{Тип данных} & \textbf{Описание} \\ \hline
\endhead
\hline
\multicolumn{2}{r}{\textit{Продолжение на следующей странице}} \\
\endfoot
\endlastfoot

% ДАННЫЕ ТАБЛИЦЫ (столько же строк, сколько и было)
bool\_ & Булев тип (True или False), хранящийся как байт \\ \hline
int\_ & Тип целочисленного значения по умолчанию (аналогичен типу long в C; обычно int64 или int32) \\ \hline
intc & Идентичен типу int в C (обычно int32 или int64) \\ \hline
intp & Целочисленное значение, используемое для индексов (аналогично типу `ssize\_t` в C; обычно int32 или int64) \\ \hline
int8 & Байт (от –128 до 127) \\ \hline
int16 & Целое число (от –32 768 до 32 767) \\ \hline
int32 & Целое число (от –2 147 483 648 до 2 147 483 647) \\ \hline
int64 & Целое число (от –9 223 372 036 854 775 808 до 9 223 372 036 854 775 807) \\ \hline
uint8 & Целое число без знака (от 0 до 255) \\ \hline
uint16 & Целое число без знака (от 0 до 65535) \\ \hline
uint32 & Целое число без знака (от 0 до 4 294 967 295) \\ \hline
uint64 & Целое число без знака (от 0 до 18 446 744 073 709 551 615) \\ \hline
float\_ & Сокращение для названия типа `float64` \\ \hline
float16 & Число с плавающей точкой с половинной точностью: 1 бит знака, 5 битов порядка, 10 битов мантиссы \\ \hline
float32 & Число с плавающей точкой с одинарной точностью: 1 бит знака, 8 битов порядка, 23 бита мантиссы \\ \hline
float64 & Число с плавающей точкой с удвоенной точностью: 1 бит знака, 11 битов порядка, 52 бита мантиссы \\ \hline
complex\_ & Сокращение для названия типа `complex128` \\ \hline
complex64 & Комплексное число, представленное двумя 32-битными числами с плавающей точкой \\ \hline
complex128 & Комплексное число, представленное двумя 64-битными числами с плавающей точкой \\ \hline
\end{longtable}

Перейдём к простейшим операциям над массивами NumPy:

\begin{itemize}
    \item получение атрибутов массивов — определение размера, формы, объема занимаемой памяти и типов 
    данных массивов:

    У каждого массива есть атрибуты \textbf{ndim} (число размерностей), \textbf{shape} (размер каждой
    размерности) и \textbf{size} (общий размер массива):

    \item индексация массивов — получение и изменение значений отдельных элементов массивов;
    \item получение срезов массивов — получение и изменение значений подмассивов внутри большого массива;
    \item изменение формы массивов — изменение формы заданного массива;
    \item объединение и разбиение массивов — объединение нескольких массивов в один и разделение одного массива на несколько.
\end{itemize}

\subsection{Библиотека Pandas}

Логичным продолжением будет рассказать про библиотеку Pandas, т.к. данная библиотека построена поверх NumPy. 
Pandas (название происходит от эконометрического термина "PANel DAta" - панельные данные, т.е. многомерные 
структурированные наборы информации) - это библиотека с открытым исходным кодом для языка Python, 
которая предоставляет высокопроизводительные и удобные инструменты для работы с данными. 

Pandas — это хорошая, а главное — масштабируемая замена привычному многим Excel. В отличие от табличного 
процессора, Pandas позволяет обрабатывать файлы объёмом в миллионы строк и автоматизировать процесс 
очистки и преобразования данных.

На практике редко приходится работать с идеально чистыми данными. Чаще в руки попадают "грязные" с 
пропусками, дубликатами, выбросами (сильно отколоняющиеся значение от общей массы значений).
При работе с данными Pandas позволяет:

\begin{enumerate}
    \item Читать данные в разных форматах: csv, txt, SQL-базы данных, JSON, HTML-страницы и т.д.
    
    Чаще всего работают с csv-файлами, выгрузками из баз данных или JSON-файлами.
    \item Чистить данные, находить пропущенные значения, убирать дубликаты, исправлять ошибки в формате дат или текста.
    \item Менять структуру таблиц, создавать новые, в общем придавать данным ту стуктуру, которая нам требуется. 
\end{enumerate}

Основными структурами данных в Pandas являются:

\begin{itemize}
    \item \textbf{Series} (Серия) — это одномерный массив данных (по сути, одна колонка таблицы).
    \item \textbf{DataFrame} (Датафрейм) — это двумерная табличная структура данных (строки и столбцы).
    При загрузке таблицы из Excel в Python, она превращается именно в DataFrame.
\end{itemize}


Установка Pandas происходит по команде в терминале (VS Code, PyCharm): 

\textbf{pip install pandas}. 
Импортируем пакет Pandas под псевдонимом \textbf{pd} (это уже в рабочем файле .py):
\begin{minted}{python}
import pandas as pd  
\end{minted}
Дальше в примерах кода буду опускать строку с импортом.

\textbf{\underline{Важное замечание!}}

Не просто так повествование начиналось с библиотеки NumPy. Для ускорения манипуляций с данными будет использоваться
векторизация. \textbf{Векторизация — это способность выполнять операции над всем массивом 
данных сразу, без явного написания цикла}.
Если приходится писать цикл \textit{for} для обработки данных в таблице, скорее всего, вы делаете что-то 
не так. Ищите векторное решение.

Ещё одно важное и удобное отличие Pandas от обычных списков или массивов NumPy — это индексы, но не 
в привычном нам понимании.

\subsection{Библиотека Scikit-learn}

\subsection{Библиотека Matplotlib / Seaborn}

\conclusion

% Отобразить все источники. Даже те, на которые нет ссылок.
% \nocite{*}

\bibliographystyle{ugost2003}
\bibliography{thesis}

% Окончание основного документа и начало приложений Каждая последующая секция
% документа будет являться приложением
\appendix

\end{document}
